\documentclass[10pt,aspectratio=169]{beamer}
\usepackage[utf8]{inputenc}
\usepackage{listings}
\usepackage{hyperref}
\usepackage{xcolor}
\usepackage{ragged2e}
\usepackage{graphicx}
\usepackage{tikz}
\usetikzlibrary{positioning,calc,decorations.pathreplacing}

\usetheme{Dresden}
%\usetheme{Bergen}

\title{Programming Techniques 2026--2027}
\author{Hannu Parviainen}
\institute{Universidad de la Laguna}
\date{\today}

\lstset{
  language=fortran,
  basicstyle=\ttfamily\scriptsize,
  columns=fullflexible,
  keepspaces=true,
  backgroundcolor=\color{white},
  keywordstyle=\color{blue},
  commentstyle=\color{gray},
  stringstyle=\color{red},
  showstringspaces=false,
  breaklines=true,
  frame=lines
}

\subtitle{Lecture 4: Data Type Precision, Programs, Modules, and Procedures}

\begin{document}

\begin{frame}
  \titlepage
\end{frame}

\input{l04a_type_precision}

\section{Program Units}

\begin{frame}
  \centering
  \Huge \textbf{Program Units}
\end{frame}


\begin{frame}{Fortran program units}
  \begin{block}{Program units}
    \begin{itemize}
      \item \textbf{program}: where execution begins. Every executable has exactly one.
      \item \textbf{module}: a container for constants, variables, types, and procedures. Other program units get access to it with \texttt{use}.
    \end{itemize}
  \end{block}
  \begin{block}{Procedures}
    \begin{itemize}
      \item \textbf{subroutine}: performs a task. It is invoked with \texttt{call} and returns no value.
      \item \textbf{function}: returns a value, so it is used inside an expression.
    \end{itemize}
  \end{block}
  \begin{block}{Where procedures live}
    Procedures are placed inside a program or a module, after a \texttt{contains} statement.
  \end{block}
\end{frame}


\begin{frame}[fragile]{The main program}
  \begin{columns}[T]
    \begin{column}{0.45\textwidth}
      \begin{itemize}
        \item Execution starts at the first executable statement of the program.
        \item The order is fixed: \texttt{use} statements, \texttt{implicit none}, declarations, and then executable statements.
        \item An optional \texttt{contains} section holds \emph{internal procedures}, which only this program can call.
        \item Here \texttt{negative} is an internal function. We return to functions later in this lecture.
      \end{itemize}
    \end{column}

    \begin{column}{0.55\textwidth}
      \begin{lstlisting}
program main
  implicit none
  real :: x                  ! declarations
  read *, x                  ! executable statements
  print *, floor(x)          ! intrinsic function
  print *, negative(x)       ! internal function
contains
  real function negative(a)
    real, intent(in) :: a
    negative = -a
  end function negative
end program main
      \end{lstlisting}
      \begin{lstlisting}[language=bash]
$ echo 2.7 | ./a.out
           2
  -2.70000005
      \end{lstlisting}
    \end{column}
  \end{columns}
\end{frame}


\section{Modules}

\begin{frame}
  \centering
  \Huge \textbf{Modules}
\end{frame}


\begin{frame}[fragile]{A first module: kinds}
  \begin{columns}[T]
    \begin{column}{0.45\textwidth}
      \begin{itemize}
        \item A module is a named container. The simplest ones hold only constants.
        \item A \texttt{use} statement makes the contents of a module available. It goes before \texttt{implicit none}.
        \item This \texttt{kinds} module does what the precision section recommended: it defines the kind once, in one place.
        \item Every part of the program then uses \texttt{real(dp)}, and changing the precision of the whole program means changing one line.
      \end{itemize}
    \end{column}

    \begin{column}{0.55\textwidth}
      \begin{lstlisting}
module kinds
  use iso_fortran_env, only: real32, real64
  implicit none
  integer, parameter :: sp = real32
  integer, parameter :: dp = real64
end module kinds

program use_kinds
  use kinds, only: dp
  implicit none
  real(dp) :: x = 0.1_dp
  print *, x   ! 0.10000000000000001
end program use_kinds
      \end{lstlisting}
    \end{column}
  \end{columns}
\end{frame}


\begin{frame}[fragile]{Modules with procedures}
  \begin{columns}[T]
    \begin{column}{0.45\textwidth}
      \begin{itemize}
        \item Most modules are libraries: they collect related constants, types, functions, and subroutines.
        \item The procedures go after \texttt{contains}.
        \item A module can share a file with the program, as here, but it must come first. Usually each module has its own file.
        \item Modules are compiled separately, so many programs can use the same one.
      \end{itemize}
    \end{column}

    \begin{column}{0.55\textwidth}
      \begin{lstlisting}
module mymath
  implicit none
  real, parameter :: pi = 3.14159
contains
  real function square(a)
    real, intent(in) :: a
    square = a**2
  end function square
end module mymath

program ex1
  use mymath
  implicit none
  print *, square(2.3), pi
end program ex1
      \end{lstlisting}
    \end{column}
  \end{columns}
\end{frame}


\begin{frame}[fragile]{Compiling a program that uses a module}
  \begin{columns}[T]
    \begin{column}{0.5\textwidth}
      \begin{itemize}
        \item With the module and the program in \texttt{ex1.f90}, compiling produces the executable \texttt{ex1} and a module file, \texttt{mymath.mod}.
        \item The \texttt{.mod} file describes the contents of the module, much like a header file in C. The compiler reads it whenever it meets \texttt{use mymath}.
      \end{itemize}
    \end{column}

    \begin{column}{0.5\textwidth}
      \begin{lstlisting}[language=bash]
$ gfortran -o ex1 ex1.f90
$ ls
ex1  ex1.f90  mymath.mod
$ ./ex1
   5.28999996       3.14159012
      \end{lstlisting}
    \end{column}
  \end{columns}
\end{frame}


\begin{frame}[fragile]{Modules in separate files}
  \begin{columns}[T]
    \begin{column}{0.52\textwidth}
      \begin{itemize}
        \item Now \texttt{mymath.f90} holds the module and \texttt{ex1.f90} the program.
        \item The module must be compiled first, because the program needs its \texttt{.mod} file.
        \item Either compile it separately and link the object file, or give both files in the right order.
        \item For larger projects, use a \textbf{Makefile} (covered later in the course).
      \end{itemize}
      \begin{lstlisting}[language=bash]
$ gfortran -c mymath.f90
$ gfortran -o ex1 ex1.f90 mymath.o

$ gfortran -o ex1 mymath.f90 ex1.f90
      \end{lstlisting}
    \end{column}

    \begin{column}{0.48\textwidth}
\begin{lstlisting}[language=Fortran,caption={mymath.f90}]
module mymath
  implicit none
  real, parameter :: pi = 3.14159
contains
  real function square(a)
    real, intent(in) :: a
    square = a**2
  end function square
end module mymath
\end{lstlisting}

\begin{lstlisting}[language=Fortran,caption={ex1.f90}]
program ex1
  use mymath
  implicit none
  print *, square(2.3), pi
end program ex1
\end{lstlisting}
    \end{column}
  \end{columns}
\end{frame}


\begin{frame}[fragile]{Controlling access: \texttt{private}, \texttt{public}, and \texttt{only}}
  \begin{columns}[T]
    \begin{column}{0.45\textwidth}
      \begin{itemize}
        \item Everything in a module is \texttt{public} by default.
        \item A bare \texttt{private} statement hides everything from the users of the module. A \texttt{public ::} statement then lists what the module offers.
        \item On the other side, \texttt{use shapes, only: sphere\_volume} imports only the listed names.
        \item This avoids name clashes and shows where each name comes from. We have already used it with \texttt{iso\_fortran\_env}.
      \end{itemize}
    \end{column}

    \begin{column}{0.55\textwidth}
      \begin{lstlisting}
module shapes
  implicit none
  private
  public :: sphere_volume
  real, parameter :: pi = 3.14159   ! private
contains
  real function sphere_volume(r)
    real, intent(in) :: r
    sphere_volume = 4.0 / 3.0 * pi * r**3
  end function sphere_volume
end module shapes

program ex2
  use shapes, only: sphere_volume
  implicit none
  print *, sphere_volume(2.0)   ! 33.5102959
end program ex2
      \end{lstlisting}
    \end{column}
  \end{columns}
\end{frame}


\section{Procedures}

\begin{frame}
  \centering
  \Huge \textbf{Procedures}
\end{frame}


\begin{frame}[fragile]{Subroutines}
  \begin{columns}[T]
    \begin{column}{0.45\textwidth}
      \begin{itemize}
        \item A subroutine performs a task and is invoked with \texttt{call}.
        \item It returns no value. Any results go back through its arguments.
        \item The general form:
      \end{itemize}
      \begin{lstlisting}
subroutine name(arguments)
  ! declarations of the arguments
  ! local declarations
  ! executable statements
end subroutine name
      \end{lstlisting}
    \end{column}

    \begin{column}{0.55\textwidth}
      \begin{lstlisting}
program figures
  implicit none
  call output_figures([1.0, 2.0])
contains
  subroutine output_figures(numbers)
    real, intent(in) :: numbers(:)
    print *, 'Here are the figures', numbers
  end subroutine output_figures
end program figures
      \end{lstlisting}
      \begin{lstlisting}[language=bash]
$ ./a.out
 Here are the figures   1.00000000       2.00000000
      \end{lstlisting}
    \end{column}
  \end{columns}
\end{frame}


\begin{frame}[fragile]{Functions}
  \begin{columns}[T]
    \begin{column}{0.45\textwidth}
      \begin{itemize}
        \item A function returns a value, so it is used inside an expression.
        \item There are two ways to name the returned value:
          \begin{itemize}
            \item with \texttt{result(res)}, where \texttt{res} is declared like any other variable,
            \item with the name of the function, whose type is then given in front of \texttt{function}.
          \end{itemize}
        \item Both do the same thing. We use the \texttt{result} form for the recursive functions below.
      \end{itemize}
    \end{column}

    \begin{column}{0.55\textwidth}
      \begin{lstlisting}
program functions
  implicit none
  ! Both functions return 5.0
  print *, length_a(3.0, 4.0), length_b(3.0, 4.0)
contains
  function length_a(a, b) result(res)
    real, intent(in) :: a, b
    real :: res
    res = sqrt(a**2 + b**2)
  end function length_a

  real function length_b(a, b)
    real, intent(in) :: a, b
    length_b = sqrt(a**2 + b**2)
  end function length_b
end program functions
      \end{lstlisting}
    \end{column}
  \end{columns}
\end{frame}


\begin{frame}[fragile]{Arguments and \texttt{intent}}
  \begin{columns}[T]
    \begin{column}{0.48\textwidth}
      \begin{itemize}
        \item The names in the definition of a procedure are its \emph{dummy arguments}. The call supplies the \emph{actual arguments}.
        \item Declare the intent of every dummy argument:
          \begin{itemize}
            \item \texttt{intent(in)}: only read,
            \item \texttt{intent(out)}: set by the procedure,
            \item \texttt{intent(inout)}: read and modified.
          \end{itemize}
        \item The compiler then catches mistakes, such as assigning to an \texttt{intent(in)} argument.
        \item \texttt{v(:)} is an \emph{assumed-shape} array: it takes its shape from the caller.
      \end{itemize}
    \end{column}

    \begin{column}{0.52\textwidth}
      \begin{lstlisting}
module vectors
  implicit none
contains
  subroutine scale(alpha, v)
    real, intent(in)    :: alpha
    real, intent(inout) :: v(:)
    v = alpha * v
  end subroutine scale
end module vectors

program scale_vector
  use vectors, only: scale
  implicit none
  real :: v(3) = [1.0, 2.0, 3.0]
  call scale(2.0, v)
  print *, v   ! 2.0  4.0  6.0
end program scale_vector
      \end{lstlisting}
    \end{column}
  \end{columns}
\end{frame}


\begin{frame}[fragile]{Recursive functions}
  \begin{columns}[T]
    \begin{column}{0.45\textwidth}
      \begin{itemize}
        \item A recursive function calls itself.
        \item This suits problems that are defined in terms of smaller versions of themselves, such as the factorial, or walking through a tree.
        \item It must be declared \texttt{recursive}. Fortran 2018 makes this the default, but gfortran still requires the keyword.
        \item It needs a case that ends the recursion. Otherwise it never stops.
      \end{itemize}
    \end{column}

    \begin{column}{0.55\textwidth}
      \begin{lstlisting}
program recursion
  implicit none
  print *, factorial(5), factorial(10)   ! 120  3628800
contains
  recursive function factorial(n) result(f)
    integer, intent(in) :: n
    integer :: f
    if (n <= 1) then
      f = 1
    else
      f = n * factorial(n - 1)
    end if
  end function factorial
end program recursion
      \end{lstlisting}
    \end{column}
  \end{columns}
\end{frame}


\begin{frame}[fragile]{Pure functions}
  \begin{columns}[T]
    \begin{column}{0.45\textwidth}
      \begin{itemize}
        \item A \texttt{pure} function has no side effects. It cannot change its arguments, which must all be \texttt{intent(in)}, change module variables, or do I/O.
        \item The same arguments therefore always give the same result.
        \item The compiler checks this promise, and can then run calls in any order. This makes pure functions safe to use in parallel code.
      \end{itemize}
    \end{column}

    \begin{column}{0.55\textwidth}
      \begin{lstlisting}
program purity
  implicit none
  print *, length(3.0, 4.0)   ! 5.0
contains
  pure function length(x, y) result(r)
    real, intent(in) :: x, y
    real :: r
    r = sqrt(x**2 + y**2)
  end function length
end program purity
      \end{lstlisting}
    \end{column}
  \end{columns}
\end{frame}


\begin{frame}[fragile]{Elemental functions}
  \begin{columns}[T]
    \begin{column}{0.45\textwidth}
      \begin{itemize}
        \item An \texttt{elemental} function is written for scalars, but can also be applied to arrays, element by element.
        \item This is how intrinsic functions such as \texttt{sin} and \texttt{exp} work on arrays.
        \item The actual arguments may be scalars or conformable arrays.
        \item \texttt{elemental} implies \texttt{pure}.
      \end{itemize}
    \end{column}

    \begin{column}{0.55\textwidth}
      \begin{lstlisting}
module squares
  implicit none
contains
  elemental real function square(x)
    real, intent(in) :: x
    square = x * x
  end function square
end module squares

program elemental_example
  use squares, only: square
  implicit none
  real :: x(3) = [1.0, 2.0, 3.0]
  print *, square(2.0)   ! 4.0
  print *, square(x)     ! 1.0  4.0  9.0
end program elemental_example
      \end{lstlisting}
    \end{column}
  \end{columns}
\end{frame}


\section{Exercises}

\begin{frame}{Exercises: machine epsilon}
  \begin{block}{Tasks}
    \begin{enumerate}
      \item Machine epsilon is the gap between 1.0 and the next larger real. Find it for 32-bit reals with a loop: start from \texttt{eps = 1.0} and keep halving it for as long as \mbox{\texttt{1.0 + eps/2 > 1.0}}.
      \item Do the same for 64-bit reals.
      \item Compare your two results with the intrinsic function \texttt{epsilon(x)}.
      \item How many significant decimal digits does each result correspond to?
    \end{enumerate}
  \end{block}
\end{frame}


\begin{frame}{Exercises: the order of a sum}
  \begin{block}{Tasks}
    \begin{enumerate}
      \item Compute the sum $\sum_{k=1}^{N} 1/k$ for $N = 10^7$ with 32-bit reals, adding the terms from $k = 1$ upwards.
      \item Compute the same sum again, now adding the terms from $k = N$ downwards.
      \item Compute it once more with 64-bit reals, and use that as the reference. Which 32-bit result is closer?
      \item Explain the difference. Hint: think about the size of the gap between neighbouring reals when a small term is added to a large sum.
    \end{enumerate}
  \end{block}
\end{frame}


\begin{frame}{Exercises: a quadratic equation}
  \begin{block}{Tasks}
    \begin{enumerate}
      \item Write a program that solves $a x^2 + b x + c = 0$ with the standard formula
        \[ x_{1,2} = \frac{-b \pm \sqrt{b^2 - 4ac}}{2a} \]
        using 32-bit reals.
      \item Test it with $a = 1$, $b = -10^4$, and $c = 1$. The exact roots are close to $10^4$ and $10^{-4}$. What do you get?
      \item Explain what went wrong with the smaller root.
      \item Fix it without changing the precision. Hint: the product of the two roots is $c/a$.
      \item Repeat with 64-bit reals. How many digits of the smaller root does the standard formula get right now?
    \end{enumerate}
  \end{block}
\end{frame}

\begin{frame}{Exercises: modules, subroutines, and functions}
  \begin{block}{Tasks}
    \begin{enumerate}
      \item Write a module \texttt{mymath} in a file \texttt{mymath.f90} that contains:
        \begin{itemize}
          \item a parameter \texttt{pi = 3.14159},
          \item a function \texttt{circle\_area(r)} that returns the area of a circle.
        \end{itemize}
      \item Write a main program in a file \texttt{test\_mymath.f90} that uses the module and prints the area of a circle with radius 2.0.
      \item Compile and link the module and the program.
      \item Add a \texttt{kinds} module that defines \texttt{dp}, and change \texttt{mymath} and the program to use \texttt{real(dp)}. How many digits of the result can you trust now?
    \end{enumerate}
  \end{block}
\end{frame}


\begin{frame}{Exercises: a subroutine with several results}
  \begin{block}{Tasks}
    \begin{enumerate}
      \item Write a module \texttt{stats} with a subroutine \texttt{mean\_and\_std(x, mean, std)} that returns the mean and the sample standard deviation of a real array \texttt{x} of any size.
      \item Give every dummy argument the right \texttt{intent}, and make the array assumed-shape.
      \item Write a main program that calls the subroutine with the magnitudes \mbox{\texttt{[4.2, 3.9, 4.5, 4.1, 4.3]}} and prints the results.
      \item Add a statement to the subroutine that changes \texttt{x}. What does the compiler say?
      \item Make the module \texttt{private}, export only the subroutine, and import it with \texttt{only}.
    \end{enumerate}
  \end{block}
\end{frame}


\begin{frame}{Exercises: an elemental function}
  \begin{block}{Tasks}
    \begin{enumerate}
      \item Write a module \texttt{photometry} with an elemental function \texttt{flux\_to\_mag(flux)} that returns the magnitude $m = -2.5 \log_{10} F$.
      \item Call it from a main program with a single flux, and then with an array of fluxes.
      \item Add a \texttt{print} statement inside the function. What does the compiler say, and why?
      \item What happens if one of the fluxes is zero or negative? Protect the calculation in the main program with a \texttt{where} construct.
    \end{enumerate}
  \end{block}
\end{frame}


\begin{frame}{Exercises: recursive functions}
  \begin{block}{Tasks}
    \begin{enumerate}
      \item Write a recursive function \texttt{factorial(n)} that computes $n!$.
      \item Test the function by computing $5!$ and $10!$.
      \item Modify the function to handle invalid inputs (e.g.\ $n < 0$).
      \item Write a recursive function \texttt{fibonacci(n)} that returns the $n$th Fibonacci number.
      \item Compare the recursive implementation with an iterative one.
    \end{enumerate}
  \end{block}
\end{frame}


\end{document}
